\clearpage
% C052 / 03_proof.tex
% 四种严格证明：最快 / 最低门槛 / 最直觉 / 常规联立
% 图形构造见 C052_definition_pythagoras.tex。
\begin{proofbox}

\textbf{\textcolor{CProof}{共同约定（只设置一次，四法共用）}}

取 $e=c/a\in(0,1)$，$c^2=a^2-b^2$。下图约定四种证明中的点、长度和方向：

\begin{center}% Exact model: a=3, b^2=5, c=2, theta=60 degrees;
% r=5/4, s=5/2. Coordinates use one uniform scale.
\begin{tikzpicture}[font=\small,>=stealth]
\draw[rounded corners=3pt,gray!45] (0,0) rectangle (14.2,10.8);
\node at (7.1,10.2) {四法共用一张图：先认点，再认两段，最后读坐标};
\begin{scope}[shift={(4.6,5.6)},scale=1.15]
\draw[->,gray] (-3.35,0)--(3.5,0) node[above] {$x$};
\draw[->,gray] (0,-2.65)--(0,2.65) node[left] {$y$};
\draw[teal!75!black,thick] (0,0) ellipse (3 and 2.23606798);
% Projections: the same color identifies each endpoint and its coordinates.
\draw[blue!55!black,densely dashed] (2.625,0)--(2.625,1.08253175)--(0,1.08253175);
\draw[orange!75!black,densely dashed] (.75,0)--(.75,-2.16506351)--(0,-2.16506351);
\draw[blue!75!black,very thick] (2,0)--(2.625,1.08253175)
 node[midway,above left=5pt] {$r$};
\draw[orange!85!black,very thick] (2,0)--(.75,-2.16506351)
 node[midway,below right=5pt] {$s$};
\fill (2.625,1.08253175) circle(2pt) node[above right=5pt] {$A$};
\fill (.75,-2.16506351) circle(2pt) node[below right=5pt] {$B$};
\fill (-2,0) circle(2pt) node[below=7pt] {$F_1(-c,0)$};
\fill (2,0) circle(2pt) node[above left=8pt] {$F=F_2$};
\node[below left=4pt] at (0,0) {$O$};
\draw[violet,->] (2.45,0) arc[start angle=0,end angle=60,radius=.45];
\node[violet] at (2.75,.35) {$\theta$};
\end{scope}
\node[anchor=west,align=left] at (9.05,6.4)
 {$F=(c,0)$\\[11pt]
  $\textcolor{blue!75!black}{r=FA>0}$\\[8pt]
  $\textcolor{orange!85!black}{s=FB>0}$\\[11pt]
  $u=\cos\theta$\\[8pt]
  正方向：$F\to A$\\[5pt]
  单位方向\\[3pt]
  $(\cos\theta,\sin\theta)$};
\node at (7.1,2.05) {虚线读投影；从 $F(c,0)$ 出发，$A$ 沿正方向，$B$ 沿反方向。};
\node at (7.1,1.25)
 {$\textcolor{blue!75!black}{A=(c+ru,\ r\sin\theta)},\qquad
   \textcolor{orange!85!black}{B=(c-su,\ -s\sin\theta)}$};
\node at (7.1,.5) {整条弦：$L=AB=r+s$；图取 $a=3,\ b^2=5,\ c=2,\ \theta=60^\circ$。};
\end{tikzpicture}
\end{center}

两端点位于焦点两侧，故 $r,s>0$；取 $u=\cos\theta$，目标是求 $L=r+s$。
图中取第一象限方向便于认字母；一般推导仍允许任意方向，投影按正负号计算。
$a>c$ 保证 $a\pm cu>0$。

\end{proofbox}
\clearpage
\begin{proofbox}
\textbf{\textcolor{CProof}{证明一：最快证明（已会椭圆焦半径时，推荐考场使用）}}

\textbf{已知工具：}对椭圆上的点 $P(x,y)$，到右焦点的距离为
$PF_2=a-ex$，其中 $e=c/a$（下文证明二、三均不要求事先记住它）。

\textbf{第1步 两段各列一个一次方程：}
\[
\begin{aligned}
 r&=a-e(c+ru)
   &&\Longrightarrow& r(a+cu)&=a^2-c^2=b^2,\\
 s&=a-e(c-su)
   &&\Longrightarrow& s(a-cu)&=a^2-c^2=b^2.
\end{aligned}
\]
\textbf{第2步 直接相加，平方差收尾：}
\[
\begin{aligned}
 L=r+s
 &=\frac{b^2}{a+cu}+\frac{b^2}{a-cu}\\
 &=\frac{b^2[(a-cu)+(a+cu)]}{(a+cu)(a-cu)}\\
 &=\boxed{\frac{2ab^2}{a^2-c^2\cos^2\theta}}.
\end{aligned}
\]
\textbf{最快记忆链：}$\displaystyle PF_2=a-ex\;\Rightarrow\;
\frac{b^2}{a+cu}+\frac{b^2}{a-cu}\;\Rightarrow\;$平方差。

\medskip
\end{proofbox}
\clearpage
\begin{proofbox}
\textbf{\textcolor{CProof}{证明二：最低门槛证明（只用椭圆定义＋勾股定理，无需预记焦半径或韦达）}}


\begin{center}\begin{tikzpicture}[font=\small,>=stealth]
\draw[rounded corners=3pt,gray!45] (0,0) rectangle (14.2,10.5);
\node at (7.1,9.95) {证明二的构造：椭圆定义给斜边，勾股定理给平方};
\begin{scope}[shift={(4.4,5.4)},scale=1.5]
\draw[->,gray] (-2.35,0)--(2.15,0) node[above] {$x$};
\draw[teal!75!black,thick] (0,0) ellipse (2 and 1.7320508);
\fill[violet!10] (-1,0)--(1.6,0)--(1.6,1.0392305)--cycle;
\draw[violet,thick] (-1,0)--(1.6,1.0392305);
\draw[blue!75!black,very thick] (1,0)--(1.6,1.0392305);
\draw[orange!85!black,thick,dashed] (-1,0)--(1.6,0)--(1.6,1.0392305);
\draw[orange!85!black] (1.4,0)--(1.4,.2)--(1.6,.2);
\fill (-1,0) circle(2pt) node[below=4pt] {$F_1$};
\fill (1,0) circle(2pt) node[below=4pt] {$F$};
\fill (1.6,0) circle(2pt) node[below=10pt] {$H$};
\fill (1.6,1.0392305) circle(2pt) node[above right=4pt] {$A$};
\draw[gray,dotted] (-1,-.7)--(-1,-2.25);
\draw[gray,dotted] (1.6,-.9)--(1.6,-2.25);
\draw[<->,orange!85!black] (-1,-2.2)--(1.6,-2.2) node[midway,below] {$F_1H=2c+r\cos\theta$};
\end{scope}
\node[anchor=west,align=left] at (9.15,6.5) {$AF=r$\\[8pt]$AF_1=2a-r$\\[8pt]$AH=r\sin\theta$\\[10pt]$AF_1^2=F_1H^2+AH^2$};
\node at (7.1,.65) {$\displaystyle (2a-r)^2=(2c+r\cos\theta)^2+(r\sin\theta)^2$};
\end{tikzpicture}
\end{center}
图取 $a=2,b^2=3,c=1,\theta=60^\circ$。一般推导只用图中的水平、竖直分量；即使投影方向改变，平方后的距离表达式仍成立。


\textbf{第1步 定义与勾股给出同一个平方：}
$AF_1=2a-r$，故
\[
 (2a-r)^2=(2c+ru)^2+r^2\sin^2\theta=4c^2+4cru+r^2.
\]
\textbf{第2步 两边消去 $r^2$：}
\[
 4a^2-4ar=4c^2+4cru
 \ \Longrightarrow\ \boxed{r(a+cu)=b^2}.
\]
\end{proofbox}
\clearpage
\begin{proofbox}
\textbf{\textcolor{CProof}{证明二（续）：补齐另一段 $BF=s$}}

\begin{center}% Same ellipse, direction and scale as C052_definition_pythagoras.tex.
% a=2, b^2=3, c=1, theta=60 degrees: B=(0,-sqrt(3)), s=2.
\begin{tikzpicture}[font=\small,>=stealth]
\draw[rounded corners=3pt,gray!45] (0,0) rectangle (14.2,10.5);
\node at (7.1,9.95) {证明二的另一半：沿反方向到 $B$，水平分量由加变减};
\begin{scope}[shift={(4.4,5.4)},scale=1.5]
\draw[->,gray] (-2.35,0)--(2.15,0) node[above] {$x$};
\draw[teal!75!black,thick] (0,0) ellipse (2 and 1.7320508);
\fill[violet!10] (-1,0)--(0,0)--(0,-1.7320508)--cycle;
\draw[violet,thick] (-1,0)--(0,-1.7320508);
\draw[orange!85!black,very thick] (1,0)--(0,-1.7320508);
\draw[blue!65!black,thick,dashed] (-1,0)--(0,0)--(0,-1.7320508);
\draw[blue!65!black] (-.18,0)--(-.18,-.18)--(0,-.18);
\fill (-1,0) circle(2pt) node[below left=4pt] {$F_1$};
\fill (1,0) circle(2pt) node[above right=4pt] {$F$};
\fill (0,0) circle(2pt) node[above=5pt] {$K$};
\fill (0,-1.7320508) circle(2pt) node[below right=5pt] {$B$};
% The full focal distance above; the remaining horizontal leg below.
\draw[gray,dotted] (-1,.45)--(-1,2.2);
\draw[gray,dotted] (1,.45)--(1,2.2);
\draw[<->,orange!85!black] (-1,2.15)--(1,2.15)
 node[midway,above] {$F_1F=2c$};
\draw[gray,dotted] (-1,-.7)--(-1,-2.3);
\draw[gray,dotted] (0,-2.02)--(0,-2.3);
\draw[<->,blue!65!black] (-1,-2.25)--(0,-2.25)
 node[midway,below] {$F_1K=|2c-s\cos\theta|$};
\end{scope}
\node[anchor=west,align=left] at (9.15,6.1)
 {$\textcolor{orange!85!black}{BF=s}$\\[8pt]
  $BF_1=2a-s$\\[8pt]
  $BK=|s\sin\theta|$\\[8pt]
  本图：$KF=s\cos\theta$\\[10pt]
  $BF_1^2=F_1K^2+BK^2$};
\node at (7.1,.65)
 {$\displaystyle (2a-s)^2=(2c-s\cos\theta)^2+(s\sin\theta)^2$};
\end{tikzpicture}
\end{center}

\textbf{第3步 看反方向的直角三角形 $BF_1K$：}
两图采用同一椭圆、方向和尺度，$K$ 是 $B$ 在 $x$ 轴上的垂足。
本图中 $K$ 位于两焦点之间，$F_1K=F_1F-KF$；一般位置用绝对值表示长度，平方后同式。
由椭圆定义 $BF_1=2a-s$，故
\[
 (2a-s)^2=(2c-su)^2+s^2\sin^2\theta
 =4c^2-4csu+s^2.
\]
消去 $s^2$，得 $\boxed{s(a-cu)=b^2}$。

\textbf{第4步 把一条弦的两半加起来：}
\[
 L=r+s=\frac{b^2}{a+cu}+\frac{b^2}{a-cu}
 =\boxed{\frac{2ab^2}{a^2-c^2u^2}}.
\]
\textbf{最低门槛记忆链：}定义 $2a$ $\to$ 勾股平方 $\to$ $r^2$ 抵消
$\to$ 两段相加。完全不要求知道焦点极坐标方程。

\medskip
\end{proofbox}
\clearpage
\begin{proofbox}
\textbf{\textcolor{CProof}{证明三：最直觉证明（整体看“总长”和“两侧差”，直接看出平方的来源）}}

\textbf{先现场推出焦半径，杜绝死记：}
对椭圆上一点 $P(x,y)$，设 $d_1=PF_1$、$d_2=PF_2$。
由椭圆定义 $d_1+d_2=2a$。另一方面
\[
 d_1^2-d_2^2
 =[(x+c)^2+y^2]-[(x-c)^2+y^2]=4cx.
\]
平方差得 $(d_1-d_2)\cdot2a=4cx$，所以
\[
 d_1-d_2=\frac{2cx}{a}
 \quad\Longrightarrow\quad
 \boxed{PF_2=d_2=a-\frac ca x=a-ex}.
\]
\textbf{这是全程唯一需要的辅助结论：}两焦半径\textbf{和固定、差由横坐标决定}。

\textbf{再用一个总长 $L$、一个差 $D$：}
设 $L=r+s$，$D=r-s$。对 $A,B$ 分别使用上式：
\[
 r=a-e(c+ru),\qquad s=a-e(c-su).
\]
两式相加与相减（$h=a-ec=b^2/a$）：
\[
 \boxed{L=2h-euD},\qquad
 \boxed{D=-euL}.
\]
\textbf{读懂这两句：}弦两端相对焦点的``不平衡量'' $D$
恰好是总长 $L$ 的 $-eu$ 倍；总长又受到这个差值的影响。
代入即可：
\[
\begin{aligned}
 L&=2h-eu(-euL)=2h+e^2u^2L,\\
 L(1-e^2u^2)&=2h,\\
 \boxed{L}&=\boxed{\frac{2h}{1-e^2\cos^2\theta}
 =\frac{2ab^2}{a^2-c^2\cos^2\theta}}.
\end{aligned}
\]
\textbf{直觉关键：}方向只经 $u=\cos\theta$ 进入；相互影响了\textbf{两次}，
所以必然是 $u^2=\cos^2\theta$，不会是一次方。

\medskip
\end{proofbox}
\clearpage
\begin{proofbox}
\textbf{\textcolor{CProof}{证明四：常规证明（直线参数方程＋韦达／判别式，保留原方案思路）}}

\textbf{第1步 带有向距离的直线参数方程：}
\[
 x=c+t\cos\theta,\qquad y=t\sin\theta.
\]
两个交点对应 $t_1>0>t_2$，故弦长 $L=t_1-t_2=|t_1-t_2|$。

\textbf{第2步 联立椭圆，化成标准二次式：}
\[
 \frac{(c+t\cos\theta)^2}{a^2}
 +\frac{t^2\sin^2\theta}{b^2}=1
 \quad\Longrightarrow\quad
 \alpha t^2+\beta t+\gamma=0,
\]
其中
\[
 \alpha=\frac{\cos^2\theta}{a^2}+\frac{\sin^2\theta}{b^2}>0,
 \quad \beta=\frac{2c\cos\theta}{a^2},
 \quad \gamma=-\frac{b^2}{a^2}.
\]
$\gamma<0$ 与 $\alpha>0$ 保证两个根异号，印证焦点在两交点之间。

\textbf{第3步 判别式出现一个好记的常量：}
\[
\begin{aligned}
 \Delta&=\beta^2-4\alpha\gamma\\
 &=\frac{4c^2\cos^2\theta}{a^4}
   +\frac{4b^2\cos^2\theta}{a^4}
   +\frac{4\sin^2\theta}{a^2}\\
 &=\frac4{a^2}(\cos^2\theta+\sin^2\theta)
 =\boxed{\frac4{a^2}}.
\end{aligned}
\]

\textbf{第4步 用根差算长度，再变为标准形式：}
\[
\begin{aligned}
 L=|t_1-t_2|
 &=\frac{\sqrt\Delta}{\alpha}
 =\frac{2/a}{\frac{\cos^2\theta}{a^2}+\frac{\sin^2\theta}{b^2}}\\
 &=\frac{2ab^2}{b^2\cos^2\theta+a^2\sin^2\theta}\\
 &=\boxed{\frac{2ab^2}{a^2-c^2\cos^2\theta}}.
\end{aligned}
\]
此法对 $\theta=90^\circ$ 也适用，\textbf{无需}斜率 $k=\tan\theta$，所以不会遗漏竖直线。

\medskip
\textbf{\textcolor{CProof}{两个无需背公式的当场核验}}

\begin{itemize}[leftmargin=*,itemsep=0.4em,topsep=0.35em]
 \item \textbf{竖直通径：}令 $x=c$，由椭圆方程有
 $y^2=b^2(1-c^2/a^2)=b^4/a^2$，上下端点分别为
 $\bigl(c,\pm b^2/a\bigr)$，所以 $L=2b^2/a$。
 \item \textbf{横向长轴弦：}焦点弦与 $x$ 轴重合，两端是 $(-a,0),(a,0)$，所以 $L=2a$。
\end{itemize}

\textbf{\textcolor{CProof}{关于左焦点：}}
以上为右焦点 $F_2$ 的证明；左焦点只需关于 $y$ 轴做镜像，
$\cos\theta$ 变号而 $\cos^2\theta$ 不变，所以\textbf{完全同式}。

\end{proofbox}
